% =============================================================
%  CS102-F26 Digital Design and Computer Architecture I
%  Lecture 05 -- Karnaugh Maps (Unit 5)
%  Follows Section 2.3.2 of A Student's Computer Architecture Book I.
%
%  Build the slides and the print handout:  npm run materials -- lecture-05
% =============================================================
\documentclass{wosm-slides}
\usepackage{cs102-lectures}[2026/10/19]

\title[Lecture 05 \textperiodcentered\ Karnaugh Maps]{Karnaugh Maps}
\subtitle{Lecture 05: simplifying Boolean functions}

% ---------- Helpers for this lecture --------------------------
% Proof steps: an equation and the law that justifies it.
\newenvironment{derivation}
  {\begin{tabular}{@{}r@{\;}l@{\qquad}l@{}}}
  {\end{tabular}}
\newcommand*\by[1]{\wosmlabel{#1}}

% Empty maps with their minterm numbers.
% Three variables: rows xy, column z, as in the book.
\newenvironment{kmapthree}
  {\begin{kmap}{4}{2}{xy}{z}{00,01,11,10}{0,1}}
  {\end{kmap}}
\newcommand*\threeindices{%
  \kindices{0}{0,1}\kindices{1}{2,3}\kindices{2}{6,7}\kindices{3}{4,5}}
% Four variables: rows wx, columns yz.
\newenvironment{kmapfour}
  {\begin{kmap}{4}{4}{wx}{yz}{00,01,11,10}{00,01,11,10}}
  {\end{kmap}}
\newcommand*\fourindices{%
  \kindices{0}{0,1,3,2}\kindices{1}{4,5,7,6}%
  \kindices{2}{12,13,15,14}\kindices{3}{8,9,11,10}}

% Group colours: the second and third group, and a group we do not need.
\colorlet{group2}{wosm-brick}
\colorlet{group3}{wosm-text}
\tikzset{spare/.style={wosm-muted, dashed}}

% =============================================================
\begin{document}

\maketitle

\begin{frame}{Today}
  \tableofcontents
\end{frame}

% =============================================================
\section{Why simplify?}

\begin{frame}{Fewer gates}
  The circuit below computes $f(x) = x'''$:
  \begin{center}
    \begin{tikzpicture}[gates]
      \node (x) at (0,0) {$x$};
      \node[not gate] (n1) at (1.2,0) {};
      \node[not gate] (n2) at (2.6,0) {};
      \node[not gate] (n3) at (4.0,0) {};
      \node (f) at (5.3,0) {$f$};
      \draw[wire] (x) -- (n1.input) (n1.output) -- (n2.input)
        (n2.output) -- (n3.input) (n3.output) -- (f);
    \end{tikzpicture}
  \end{center}

  Since $x'' = x$ (involution), $f(x) = x'$: one NOT gate does the same job.

  \medskip
  \begin{itemize}
    \item Both circuits give the same output, but the long one needs three gates
      instead of one: it costs more, and the signal is slower to pass through it.
    \item Fewer and smaller gates mean cheaper and faster circuits. Minimising
      Boolean functions is worth the effort.
  \end{itemize}
\end{frame}

\begin{frame}{Canonical forms are long}
  \begin{itemize}
    \item A \alert{literal} is a variable or its complement: $x$, $x'$, \dots
    \item From Lecture 04:
      $\sum m(1, 3, 5, 6, 7) = x'y'z + x'yz + xy'z + xyz' + xyz$ has 5 products and
      15 literals.
    \item The same function is $xy + z$: 2 products and 3 literals.
  \end{itemize}

  \medskip
  The laws of Lecture 03 get you there, but you have to see which law to apply
  next. A \alert{Karnaugh map} (K-map) turns the search into a picture.

  \medskip
  \wosmlabel{Introduced by Maurice Karnaugh in 1953.}
\end{frame}

% =============================================================
\section{The Karnaugh map}

\begin{frame}{Neighbours combine}
  Two minterms that differ in exactly one variable combine into one shorter
  product:

  \medskip
  \small
  \begin{derivation}
    $xyz' + xyz$ & $= xy(z' + z) = xy \cdot 1 = xy$
      & \by{distributivity, complement, identity} \\
  \end{derivation}

  \medskip
  \normalsize
  And two such pairs, differing in one more variable, combine again:
  \[ x'y'z + x'yz + xy'z + xyz = x'z + xz = z . \]

  \begin{block}{The idea}
    Draw the truth table as a grid in which cells next to each other differ in
    exactly one input bit. Then the products that combine are the cells that sit
    together.
  \end{block}
\end{frame}

\begin{frame}{Gray code}
  To order the rows so that neighbours differ in one bit, do not count in binary:
  \begin{center}
    \begin{tabular}{@{}lccccc@{}}
      binary order & 00 & 01 & 10 & 11 & \\
      bits that change & & 1 & \alert{2} & 1 & \alert{2} back to 00 \\
      \midrule
      Gray code & 00 & 01 & 11 & 10 & \\
      bits that change & & 1 & 1 & 1 & 1 back to 00 \\
    \end{tabular}
  \end{center}

  \begin{definition}[Gray code]
    A \emph{Gray code} lists the $n$-bit numbers so that each one differs from the
    next in a single bit. For two bits: 00, 01, 11, 10.
  \end{definition}

  The last code also differs from the first in one bit: the order wraps around.
\end{frame}

\begin{frame}{A map of three variables}
  \begin{columns}[c, totalwidth=\textwidth]
    \column{0.36\textwidth}
      \centering
      \begin{kmapthree}
        \kshade{1}{0}\kshade{3}{0}\kshade{2}{1}
        \threeindices
        \node[font=\small] at (0.5,-2.5) {$\star$};
      \end{kmapthree}
    \column{0.6\textwidth}
      \begin{itemize}
        \item Rows: $xy$ in Gray-code order. Columns: $z$.
        \item Each cell is one row of the truth table; the small number is its
          minterm number.
        \item The cell $\star$ is input 110. Its neighbours, shaded, are 010, 100
          and 111: each differs from 110 in one bit.
        \item The map \alert{wraps around}: the top row (00) and the bottom row
          (10) are neighbours too.
      \end{itemize}
  \end{columns}
\end{frame}

\begin{frame}{Filling in the map}
  Simplify $f(x, y, z) = (x + y + z)(x + y' + z)(x' + y' + z)(x' + y' + z')$.

  \medskip
  \begin{columns}[c, totalwidth=\textwidth]
    \column{0.36\textwidth}
      \centering
      \begin{kmapthree}
        \threeindices
        \only<1>{\kmaprow{0}{0,}\kmaprow{1}{0,}\kmaprow{2}{0,0}}
        \only<2->{\kmaprow{0}{0,1}\kmaprow{1}{0,1}\kmaprow{2}{0,0}\kmaprow{3}{1,1}}
      \end{kmapthree}
    \column{0.6\textwidth}
      \begin{enumerate}
        \item The four maxterms are $M_0$, $M_2$, $M_6$ and $M_7$: write a 0 in
          those cells.
        \item<2-> Every other cell is a 1: $f = \sum m(1, 3, 4, 5)$.
      \end{enumerate}
  \end{columns}

  \medskip
  \uncover<2->{From a truth table, copy each row's output into its cell. Mind the
  order of the rows: 00, 01, \alert{11, 10}.}
\end{frame}

% =============================================================
\section{Grouping}

\begin{frame}{Groups}
  \begin{block}{Rules for groups}
    \begin{itemize}
      \item A group contains only 1s.
      \item It is a rectangle of 1, 2, 4, 8, \dots{} cells: a power of two.
      \item It may wrap around the edges of the map.
      \item Groups may overlap.
    \end{itemize}
  \end{block}

  \begin{itemize}
    \item Make every group \alert{as large as possible}, and use \alert{as few
      groups as possible} to cover all the 1s.
    \item Each group becomes one product. Keep the variables that stay the same in
      all its cells (as they are if 1, complemented if 0) and drop the others.
    \item In a map of $n$ variables, a group of $2^k$ cells has $n - k$ literals.
  \end{itemize}
\end{frame}

\begin{frame}{Reading the groups}
  \begin{columns}[c, totalwidth=\textwidth]
    \column{0.36\textwidth}
      \centering
      \begin{kmapthree}
        \threeindices
        \kmaprow{0}{0,1}\kmaprow{1}{0,1}\kmaprow{2}{0,0}\kmaprow{3}{1,1}
        \kgroup{0}{1}{1}{1}
        \kgroup[group2]{3}{0}{3}{1}
      \end{kmapthree}
    \column{0.6\textwidth}
      \begin{itemize}
        \item \textcolor{wosm-accent}{Teal}, cells 001 and 011: $x = 0$ and $z = 1$
          stay the same, $y$ changes. Product: $x'z$.
        \item \textcolor{group2}{Brick}, cells 100 and 101: $x = 1$ and $y = 0$ stay
          the same, $z$ changes. Product: $xy'$.
      \end{itemize}

      \medskip
      \[ f(x, y, z) = xy' + x'z \]
      \wosmlabel{4 products and 12 literals became 2 and 4.}
  \end{columns}
\end{frame}

\begin{frame}{Prime and essential prime implicants}
  \begin{columns}[c, totalwidth=\textwidth]
    \column{0.36\textwidth}
      \centering
      \begin{kmapthree}
        \threeindices
        \kmaprow{0}{0,1}\kmaprow{1}{0,1}\kmaprow{2}{0,0}\kmaprow{3}{1,1}
        \kgroup{0}{1}{1}{1}
        \kgroup[group2]{3}{0}{3}{1}
        \kgroup[spare]{-1}{1}{0}{1}\kgroup[spare]{3}{1}{4}{1}
      \end{kmapthree}
    \column{0.6\textwidth}
      \small
      \begin{definition}[Prime implicant]
        A group of 1s that cannot be made any larger.
      \end{definition}
      \begin{definition}[Essential prime implicant]
        A prime implicant that covers a 1 which no other prime implicant covers.
      \end{definition}
      The dashed group, cells 001 and 101 (wrapping around), is $y'z$. It is prime,
      but not essential: both its 1s are covered already.
  \end{columns}

  \medskip
  \wosmlabel{Coming back: $xy' + x'z + y'z = xy' + x'z$ is the consensus theorem of
  Lecture 03, and $y'z$ is the consensus term.}
\end{frame}

\begin{frame}{The procedure}
  \begin{enumerate}
    \item Fill in the map from the truth table or the canonical form.
    \item Find the prime implicants: the largest groups of 1s.
    \item Take every essential prime implicant.
    \item If any 1s are still uncovered, cover them with as few further prime
      implicants as possible.
    \item Write one product per group, and add them up.
  \end{enumerate}

  \bigskip
  \begin{remark}
    The result is a minimal sum of products. It need not be unique: sometimes two
    different choices in step 4 are equally small.
  \end{remark}
\end{frame}

\begin{frame}{Back to Lecture 04}
  \begin{columns}[c, totalwidth=\textwidth]
    \column{0.36\textwidth}
      \centering
      \begin{kmapthree}
        \threeindices
        \kmaprow{0}{0,1}\kmaprow{1}{0,1}\kmaprow{2}{1,1}\kmaprow{3}{0,1}
        \kgroup{0}{1}{3}{1}
        \kgroup[group2]{2}{0}{2}{1}
      \end{kmapthree}
    \column{0.6\textwidth}
      $f = \sum m(1, 3, 5, 6, 7)$.
      \begin{itemize}
        \item \textcolor{wosm-accent}{Teal}: the whole column $z = 1$, four cells.
          Only $z$ stays the same: $z$.
        \item \textcolor{group2}{Brick}: row 11, two cells: $xy$. It overlaps the
          teal group, and that is fine.
      \end{itemize}
      \medskip
      \[ f(x, y, z) = xy + z \]
      \wosmlabel{The expression we started from in Lecture 04.}
  \end{columns}
\end{frame}

% =============================================================
\section{Four variables}

\begin{frame}{A map of four variables}
  \begin{columns}[c, totalwidth=\textwidth]
    \column{0.42\textwidth}
      \centering
      \begin{kmapfour}
        \fourindices
      \end{kmapfour}
    \column{0.54\textwidth}
      \begin{itemize}
        \item Rows: $wx$; columns: $yz$; both in Gray-code order.
        \item Cell number $= 4 \times$ row value $+$ column value: row 11, column 01
          is $m_{13}$.
        \item The map wraps around both ways: top and bottom rows are neighbours,
          and so are the left and right columns.
        \item So the four \alert{corners} are neighbours, and can form one group.
      \end{itemize}
  \end{columns}
\end{frame}

\begin{frame}{Example: wrap-around}
  \begin{columns}[c, totalwidth=\textwidth]
    \column{0.42\textwidth}
      \centering
      \begin{kmapfour}
        \fourindices
        \kmaprow{0}{0,1,1,0}\kmaprow{1}{1,1,0,0}
        \kmaprow{2}{1,1,0,0}\kmaprow{3}{0,1,1,0}
        \kgroup{-1}{1}{0}{2}\kgroup{3}{1}{4}{2}
        \kgroup[group2]{1}{0}{2}{1}
        \only<2->{\kgroup[spare]{0}{1}{3}{1}}
      \end{kmapfour}
    \column{0.54\textwidth}
      $f = \sum m(1, 3, 4, 5, 9, 11, 12, 13)$.
      \begin{itemize}
        \item \textcolor{wosm-accent}{Teal}, wrapping top to bottom: $x = 0$,
          $z = 1$. Product: $x'z$.
        \item \textcolor{group2}{Brick}: $x = 1$, $y = 0$. Product: $xy'$.
        \item<2-> Dashed: the column $yz = 01$ is $y'z$, prime but not essential.
      \end{itemize}
      \medskip
      \[ f(w, x, y, z) = x'z + xy' \]
      \wosmlabel{Each group has 4 cells: $4 - 2 = 2$ literals.}
  \end{columns}
\end{frame}

\begin{frame}{Example: the four corners}
  \begin{columns}[c, totalwidth=\textwidth]
    \column{0.42\textwidth}
      \centering
      \begin{kmapfour}
        \fourindices
        \kmaprow{0}{1,0,0,1}\kmaprow{1}{0,1,1,0}
        \kmaprow{2}{0,1,1,0}\kmaprow{3}{1,0,0,1}
        \kgroup{-1}{-1}{0}{0}\kgroup{-1}{3}{0}{4}
        \kgroup{3}{-1}{4}{0}\kgroup{3}{3}{4}{4}
        \kgroup[group2]{1}{1}{2}{2}
      \end{kmapfour}
    \column{0.54\textwidth}
      $f = \sum m(0, 2, 5, 7, 8, 10, 13, 15)$.
      \begin{itemize}
        \item \textcolor{wosm-accent}{Teal}, the four corners: $x = 0$, $z = 0$.
          Product: $x'z'$.
        \item \textcolor{group2}{Brick}, the centre: $x = 1$, $z = 1$. Product:
          $xz$.
      \end{itemize}
      \medskip
      \[ f(w, x, y, z) = x'z' + xz \]
      \wosmlabel{$f$ is 1 when $x$ and $z$ are equal; $w$ and $y$ do not matter.}
  \end{columns}
\end{frame}

% =============================================================
\section{Don't cares and products of sums}

\begin{frame}{Don't-care conditions}
  Sometimes some inputs never occur, or their output does not matter. Mark them X:
  each X may be grouped as a 1 or left out, whichever makes the groups larger.

  \medskip
  \begin{columns}[c, totalwidth=\textwidth]
    \column{0.42\textwidth}
      \centering
      \begin{kmapfour}
        \fourindices
        \kmaprow{0}{0,0,0,0}\kmaprow{1}{0,1,1,1}
        \kmaprow{2}{X,X,X,X}\kmaprow{3}{1,1,X,X}
        \kgroup[group3]{2}{0}{3}{3}
        \kgroup{1}{1}{2}{2}
        \kgroup[group2]{1}{2}{2}{3}
      \end{kmapfour}
    \column{0.54\textwidth}
      \small
      The input $wxyz$ is a decimal digit, 0 to 9, in binary; inputs 10 to 15 never
      occur. $f$ is 1 when the digit is 5 or more.
      \begin{itemize}
        \item With the X's: $f = w + xz + xy$.
        \item Without them: $f = w'xz + w'xy + wx'y'$, three products of three
          literals.
      \end{itemize}
  \end{columns}
\end{frame}

\begin{frame}{Products of sums from a map}
  \begin{columns}[c, totalwidth=\textwidth]
    \column{0.36\textwidth}
      \centering
      \begin{kmapthree}
        \threeindices
        \kmaprow{0}{0,1}\kmaprow{1}{0,1}\kmaprow{2}{0,0}\kmaprow{3}{1,1}
        \kgroup{0}{0}{1}{0}
        \kgroup[group2]{2}{0}{2}{1}
      \end{kmapthree}
    \column{0.6\textwidth}
      Group the \alert{0s} instead. Each group becomes one sum: keep the variables
      that stay the same, as they are if 0 and complemented if 1.
      \begin{itemize}
        \item \textcolor{wosm-accent}{Teal}, cells 000 and 010: $x = 0$, $z = 0$.
          Sum: $x + z$.
        \item \textcolor{group2}{Brick}, cells 110 and 111: $x = 1$, $y = 1$. Sum:
          $x' + y'$.
      \end{itemize}
      \[ f = (x + z)(x' + y') \]
  \end{columns}

  \medskip
  \wosmlabel{Multiply out: $xy' + x'z + y'z$, and the consensus term $y'z$ drops: the
  same $f = xy' + x'z$ as before.}
\end{frame}

\begin{frame}{Try it}
  \begin{columns}[c, totalwidth=\textwidth]
    \column{0.42\textwidth}
      \centering
      \begin{kmapfour}
        \fourindices
        \kmaprow{0}{1,1,0,1}\kmaprow{1}{0,1,0,0}
        \kmaprow{2}{0,0,0,0}\kmaprow{3}{1,1,0,1}
        \only<2->{%
          \kgroup{-1}{-1}{0}{0}\kgroup{-1}{3}{0}{4}
          \kgroup{3}{-1}{4}{0}\kgroup{3}{3}{4}{4}
          \kgroup[group2]{-1}{0}{0}{1}\kgroup[group2]{3}{0}{4}{1}
          \kgroup[group3]{0}{1}{1}{1}}
      \end{kmapfour}
    \column{0.54\textwidth}
      \begin{enumerate}
        \item Minimise $f = \sum m(0, 1, 2, 5, 8, 9, 10)$.\par
          \uncover<2->{\small $f = x'z' + x'y' + w'y'z$: the corners, a wrapped
          group of four, and a pair.}
        \item Minimise the majority function of Lecture 04, $\sum m(3, 5, 6, 7)$.\par
          \uncover<3->{\small $xy + xz + yz$: three pairs, all essential.}
      \end{enumerate}
  \end{columns}
\end{frame}

% =============================================================
\section{Summary}

\begin{frame}{Summary}
  \begin{itemize}
    \item Minimal expressions mean fewer gates: cheaper and faster circuits.
    \item A K-map is a truth table arranged in Gray-code order, so that neighbouring
      cells differ in one bit. It wraps around at the edges.
    \item Group the 1s in rectangles of $2^k$ cells, as large and as few as
      possible; each group is one product of $n - k$ literals.
    \item Essential prime implicants are always needed; others only to cover what
      is left.
    \item Don't-cares (X) may be used as 1s. Grouping the 0s gives a product of sums.
    \item K-maps work well up to four variables. Beyond that, computers minimise
      functions, with methods such as Quine--McCluskey.
  \end{itemize}
\end{frame}

\begin{frame}{Before the next lecture}
  \begin{itemize}
    \item Read Section~2.3.2 of the book, \emph{Boolean Function Simplification with
      Karnaugh Maps}.
    \item Minimise $\sum m(0, 2, 4, 5, 6, 7)$ as a sum of products and as a product
      of sums.
    \item Minimise $\sum m(0, 1, 5, 7, 13, 15)$ in four variables, and find a
      prime implicant you do not need.
    \item Open Logisim-evolution and find the AND, OR and NOT gates.
  \end{itemize}

  \begin{block}{Next lecture \unitref{6}}
    Logic gates: NOT, AND, OR, XOR, NAND, NOR, XNOR, and universal gate sets.
  \end{block}
\end{frame}

\end{document}
